La oferta total disponible es:

\[
50+40+30=120
\]

y la demanda total es:

\[
45+35+50+40=170.
\]

Por tanto, existe un exceso de demanda de:

\[
170-120=50.
\]

Para equilibrar el problema se introduce un origen ficticio, $\mathcal{O}_4$, con una oferta de 50 lotes. Este origen ficticio representa la demanda no satisfecha. Los costes desde este origen ficticio hasta cada hospital son las penalizaciones por lote no entregado.

\[
\begin{array}{c|cccc|c}
 & \mathcal{D}_1 & \mathcal{D}_2 & \mathcal{D}_3 & \mathcal{D}_4 & O_i \\
\hline
\mathcal{O}_1 & 12 & 18 & 25 & 16 & 50 \\
\mathcal{O}_2 & 20 & 14 & 17 & 22 & 40 \\
\mathcal{O}_3 & 24 & 21 & 13 & 19 & 30 \\
\mathcal{O}_4 & 90 & 75 & 110 & 60 & 50 \\
\hline
D_j & 45 & 35 & 50 & 40 & \\
\end{array}
\]

donde $\mathcal{O}_4$ representa la demanda no atendida y, por tanto, los valores de la fila $\mathcal{O}_4$ no son costes de transporte, sino penalizaciones.

\subsubsection*{Solución inicial por el método de Vogel}

Aplicando el método de Vogel se obtiene la siguiente secuencia de asignaciones:

\[
\begin{array}{c|c|c|c}
\text{Paso} & \text{Fila/columna elegida} & \text{Variable asignada} & \text{Asignación} \\
\hline
1 & \mathcal{O}_4 & x_{44} & 40 \\
2 & \mathcal{O}_4 & x_{42} & 10 \\
3 & \mathcal{D}_1 & x_{11} & 45 \\
4 & \mathcal{O}_3 & x_{33} & 30 \\
5 & \mathcal{D}_3 & x_{23} & 20 \\
6 & \mathcal{D}_2 & x_{22} & 20 \\
7 & \mathcal{D}_2 & x_{12} & 5 \\
\end{array}
\]

Por tanto, la solución inicial obtenida mediante Vogel es la siguiente:

\begin{center}
\begin{tabular}{c|rrr|rrr|rrr|rrr|c}
& \multicolumn{3}{c|}{$\mathcal{D}_1$} 
& \multicolumn{3}{c|}{$\mathcal{D}_2$} 
& \multicolumn{3}{c|}{$\mathcal{D}_3$} 
& \multicolumn{3}{c|}{$\mathcal{D}_4$} 
& $O_i$ \\

\hline
& \footnotesize{12} &&& \footnotesize{18} &&& \footnotesize{25} &&& \footnotesize{16} &&& \\
&& \textbf{45} &&& \textbf{5} &&& &&& && 50 \tiny{(=45+5)} \\
\multirow{-3}{*}{$\mathcal{O}_1$} &&&&&&&&&&&&& \\

\hline
& \footnotesize{20} &&& \footnotesize{14} &&& \footnotesize{17} &&& \footnotesize{22} &&& \\
&& &&& \textbf{20} &&& \textbf{20} &&& && 40 \tiny{(=20+20)} \\
\multirow{-3}{*}{$\mathcal{O}_2$} &&&&&&&&&&&&& \\

\hline
& \footnotesize{24} &&& \footnotesize{21} &&& \footnotesize{13} &&& \footnotesize{19} &&& \\
&& &&& &&& \textbf{30} &&& && 30 \tiny{(=30)} \\
\multirow{-3}{*}{$\mathcal{O}_3$} &&&&&&&&&&&&& \\

\hline
& \footnotesize{90} &&& \footnotesize{75} &&& \footnotesize{110} &&& \footnotesize{60} &&& \\
&& &&& \textbf{10} &&& &&& \textbf{40} && 50 \tiny{(=10+40)} \\
\multirow{-3}{*}{$\mathcal{O}_4$} &&&&&&&&&&&&& \\

\hline
$D_j$ 
&& 45 \tiny{(=45)} &&
& 35 \tiny{(=5+20+10)} &&
& 50 \tiny{(=20+30)} &&
& 40 \tiny{(=40)} && \\
\end{tabular}
\end{center}

La solución tiene $4+4-1=7$ variables básicas, por lo que es una solución básica factible no degenerada.

\subsubsection*{Comprobación de optimalidad mediante stepping stone}

Las variables básicas son:

\[
x_{11}, x_{12}, x_{22}, x_{23}, x_{33}, x_{42}, x_{44}.
\]

Para comprobar si la solución es óptima, se calcula la variación del coste total asociada a introducir cada variable no básica. En un problema de minimización, si todas estas variaciones son no negativas, la solución es óptima.

\[
\begin{array}{c|c}
\text{Variable no básica} & \Delta z \\
\hline
x_{13} & 25-18+14-17 = 4 \\
x_{14} & 16-18+75-60 = 13 \\
x_{21} & 20-12+18-14 = 12 \\
x_{24} & 22-14+75-60 = 23 \\
x_{31} & 24-13+17-14+18-12 = 20 \\
x_{32} & 21-13+17-14 = 11 \\
x_{34} & 19-13+17-14+75-60 = 24 \\
x_{41} & 90-75+18-12 = 21 \\
x_{43} & 110-75+14-17 = 32 \\
\end{array}
\]

Como todas las variaciones son positivas, no existe ninguna variable no básica cuya entrada reduzca el coste total. Por tanto, la solución obtenida mediante Vogel es ya la solución óptima.

\subsubsection*{Solución óptima}

La solución óptima es:

\[
x_{11}=45, \quad x_{12}=5, \quad x_{22}=20, \quad x_{23}=20, \quad x_{33}=30,
\quad x_{42}=10, \quad x_{44}=40.
\]

Todas las demás variables toman valor cero.

En términos del problema real:

\begin{itemize}
    \item Desde el almacén $\mathcal{O}_1$ se envían 45 lotes al hospital $\mathcal{D}_1$ y 5 lotes al hospital $\mathcal{D}_2$.
    \item Desde el almacén $\mathcal{O}_2$ se envían 20 lotes al hospital $\mathcal{D}_2$ y 20 lotes al hospital $\mathcal{D}_3$.
    \item Desde el almacén $\mathcal{O}_3$ se envían 30 lotes al hospital $\mathcal{D}_3$.
    \item Quedan sin entregar 10 lotes del hospital $\mathcal{D}_2$ y 40 lotes del hospital $\mathcal{D}_4$.
\end{itemize}

La demanda no satisfecha queda recogida en las variables del origen ficticio:

\[
x_{41}=0, \qquad x_{42}=10, \qquad x_{43}=0, \qquad x_{44}=40.
\]

Por tanto:

\[
\begin{array}{c|cccc}
 & \mathcal{D}_1 & \mathcal{D}_2 & \mathcal{D}_3 & \mathcal{D}_4 \\
\hline
\text{Demanda no satisfecha} & 0 & 10 & 0 & 40 \\
\end{array}
\]

\subsubsection*{Coste total}

El coste total de transporte y penalización es:

\[
\begin{aligned}
Z ={}& 12\cdot 45 + 18\cdot 5 + 14\cdot 20 + 17\cdot 20 + 13\cdot 30 \\
&+ 75\cdot 10 + 60\cdot 40.
\end{aligned}
\]

Por tanto:

\[
Z = 540+90+280+340+390+750+2400 = 4790.
\]

La solución óptima tiene un coste total de:

\[
{z^*=4790}
\]

Este valor incluye tanto los costes reales de transporte como las penalizaciones por demanda no satisfecha.